% Replace all of Section IV-B, "Constrained Cramer--Rao Bound".
\subsection{Constrained Cram\'er--Rao Bound}
\label{sec:revised-ccrb}

The ordinary CRLB treats the real and imaginary parts of each channel
coefficient as separate unknowns. The geometric CCRB also uses the fact
that these coefficients are generated by a common set of paths. It is
\begin{equation}
\mathbf C_{\rm geom}
=\mathbf U(\mathbf U^T\mathbf J\mathbf U)^{-1}\mathbf U^T,
\label{eq:revision-ccrb}
\end{equation}
where $\mathbf J$ is the Fisher information matrix in real channel
coordinates. The columns of $\mathbf U$ form an orthonormal basis for the
local changes in the channel allowed by the geometric model
\cite{gorman1990,stoica1998}. Appendix~\ref{app:revision-ccrb} gives the
parameter ordering, derivatives and numerical calculation.

The bound is evaluated at the true channel and its path configuration.
The path gains and spatial frequencies are still unknown coordinates in
the calculation; they are not supplied to the estimator. The known
quantities are the pilots, reference field and noise variance. The true
model order is used only to define this ideal geometric reference.

We plot the trace of each bound with the same channel-energy normalization
as the NMSE. The gap between the two references shows the potential value
of the geometric information under their local unbiasedness assumptions.
It does not prove that the gap is achievable by HS-GS or another practical
estimator. EM-GS and HS-GS are not assumed to be unbiased, so neither curve
is treated as an absolute lower limit on their measured NMSE.
